\section{Introduction}
\label{sec:intro}

How people find web content has changed remarkably little since the advent of search engines.
Users type keywords, follow hyperlinks, or accept algorithmic recommendations---all of which organize the web by \emph{topic}, authority, or past behavior.
Yet when someone describes a web page informally, they rarely lead with its subject matter.
Instead, they reach for aesthetic shorthand: a site ``gives early-2010s personal blog energy,'' or ``has that corporate LinkedIn vibe.''
This subjective, cultural dimension---what we call the \emph{vibe} of a page---is invisible to every existing navigation system.

{\bf Why does vibe matter?}
Vibes capture information orthogonal to topic: a Romanian collaboration blog and a British movie review site may share nothing in subject matter, yet both project the same warm, hand-crafted, personal-blog energy.
A user looking for content that \emph{feels} a certain way has no tool to find it; current systems cannot surface the latent aesthetic structure of the web.
Recent work has begun to formalize the notion of ``vibes'' as measurable axes of variation in LLM outputs~\cite{dunlap2024vibecheck}, but this line of research has not yet been applied to web content at large.

{\bf Our approach.}
We propose a simple three-stage pipeline (\figref{fig:maps}): (1)~prompt a large language model with the text of a web page and ask ``\emph{What is this web page giving? Start your answer with `it's giving'},'' (2)~embed the resulting vibe description with a sentence transformer, and (3)~project the embeddings to two dimensions with UMAP to produce an interactive \emph{vibe map}.
We compare this vibe space against a baseline content space formed by directly embedding the raw page text.
On 497 pages from the \cfour corpus~\cite{raffel2020exploring}, vibe space and content space disagree substantially: the Adjusted Rand Index between their clusterings is only 0.113 ($p{=}0.0002$), and vibe-space neighbors exhibit a Cohen's $d{=}1.44$ reduction in raw-text similarity relative to content-space neighbors.
We find 293 cross-domain page pairs that cluster together in vibe space despite covering entirely different topics.

Our contributions are:
\begin{itemize}[leftmargin=*,itemsep=0pt,topsep=0pt]
    \item We introduce a pipeline that transforms web pages into vibe descriptions using a culturally grounded LLM prompt and show that the resulting embeddings capture aesthetic similarity distinct from topical similarity.
    \item We provide quantitative evidence that vibe space organizes pages differently from content space, with strong statistical significance and large effect sizes across multiple evaluation metrics.
    \item We demonstrate cross-domain discovery: pages about unrelated subjects cluster together when they share aesthetic character, enabling a fundamentally new mode of web browsing.
\end{itemize}

The remainder of this paper is organized as follows.
\Secref{sec:related} reviews related work on embedding visualization, LLM-based content description, and style-aware representations.
\Secref{sec:method} describes our pipeline and evaluation methodology.
\Secref{sec:results} presents quantitative and qualitative results, and \secref{sec:discussion} discusses limitations and future directions.
